% QUESTAO
Dadas duas matrizes $A$ e $B$ de mesma ordem, denominamos diferença entre $A$ e $B$, indicada por $A-B$, a matriz $C$ obtida ao calcularmos a adição de $A$ com o oposto de $B$, ou seja, $A-B=A+(-B)=C$. Assim, Calcule a matriz $C$, tal que $C=\left[ \begin{array}{cc}
10 & 18 \\
-9 & 15
\end{array} \right] - \left[ \begin{array}{cc}
-7 & 0 \\
12 & -4
\end{array} \right]$

% ALTERNATIVAS
$C=\left[ \begin{array}{cc}	17 & 18 \\	-21 & 19 \end{array} \right]$
$C=\left[ \begin{array}{cc}	3 & 18 \\	-17 & 19 \end{array} \right]$
$C=\left[ \begin{array}{cc}	20 & 0 \\	21 & 19 \end{array} \right]$
$C=\left[ \begin{array}{cc}	17 & 18 \\	-21 & 11 \end{array} \right]$
$C=\left[ \begin{array}{cc}	7 & 18 \\	-3 & 19 \end{array} \right]$


